\emph{Setting.} In a causal tower $e_{t}$ is the basis vector of the minority sign, $e_{a}$ a space direction, $Q(e_{t})Q(e_{a}) = -1$, $f = \tfrac12(1 + e_{t}e_{a})$ and $S = \Cl^{0}(V)f$, of real dimension $2$, $4$, $8$ for $n = 3, 4, 5$ (Theorem \claimref{dark-tower-matter.constituents-field}; the check rebuilds a basis of $S$ and confirms the dimensions). For a rotor $R = \exp(-B/2)$ and $\psi \in S$, $R\psi \in S$ and $R\psi\, e_{t}\rev{(R\psi)} = R(\psi e_{t}\rev{\psi})\rev{R}$; conjugation by $R$ preserves grades, so $j(R\psi) = Rj(\psi)\rev{R}$. Differentiating at $R = 1$, the symmetric bilinear map $\beta_{j}(\psi, \phi) = \tfrac12\langle \psi e_{t}\rev{\phi} + \phi e_{t}\rev{\psi}\rangle_{1}$ satisfies the covariance equation of the statement. Conversely that equation for all bivectors is covariance under the connected group the rotors generate.

(i) A symmetric bilinear map $\beta : S \times S \to V$ is given by its values $\beta(b_{i}, b_{k}) \in V$ on a basis $b_{1}, \ldots, b_{d}$ of $S$, $i \le k$, that is by $\tfrac12 d(d+1)n$ real coordinates ($9$, $40$, $180$ for $n = 3, 4, 5$). The covariance equation, written for each basis bivector $B$, each pair $i \le k$ and each component of $V$, is a homogeneous linear system in these coordinates whose coefficients are the matrix entries of $\psi \mapsto B\psi$ on the basis of $S$ and of $v \mapsto Bv - vB$ on the basis of $V$. These entries are rational, so the rank of the system over $\Q$ equals its rank over $\R$. The check computes the rank exactly in each of the six towers: the solution space has dimension $1$ in every one. By the Setting $\beta_{j}$ is a solution, and it is not zero (the check confirms both on the basis). So every spin-covariant quadratic vector is a real multiple of $j$.

(ii) Let $y = \psi e_{t}\rev{\psi}$ with $\psi$ even. Then $y$ is odd and $\rev{y} = y$. Reversion multiplies grade $k$ by $(-1)^{k(k-1)/2}$, which is $+1$ on odd grades exactly for $k \equiv 1 \pmod 4$. So $y$ has grades $1$ and $5$ only: for $n \le 4$ it is a vector, and for $n = 5$ it is $j + sI$ with $I$ the pseudoscalar. In odd dimension $I$ commutes with every vector, hence is central (the check confirms it on the generators). The check, with $\psi$ symbolic on a basis of $S$, confirms that the grade-$5$ coefficient $s$ is not the zero polynomial for $n = 5$ and that $j$ is not zero in every tower. Since $I$ is central and conjugation by a rotor fixes it, $s(R\psi) = s(\psi)$, so $s$ is spin-invariant. The space of spin-invariant quadratic forms is the solution space of the scalar-valued system, of dimension $0$ for $n = 3, 4$ and $1$ for $n = 5$ by the same exact rank computation; for $n = 5$ it is therefore the line of $s$.

(iii) Let $n = 5$. Since $\psi = \psi f$, $y = \psi\,\nu\,\rev{\psi}$ with $\nu = f e_{t}\rev{f} = \tfrac12(e_{t} - Q(e_{t})e_{a})$, a null vector: $Q(\nu) = \tfrac14(Q(e_{t}) + Q(e_{a})) = 0$. The element $\rev{\psi}\psi$ is even and fixed by reversion, so it has grades $0$ and $4$ only: $\rev{\psi}\psi = \alpha + Ib$ with $\alpha$ a scalar and $b$ a vector, because $I$ is central and invertible. Then $\nu(\alpha + Ib)\nu = \alpha Q(\nu) + I\,\nu b\nu = I\bigl(2(\nu\cdot b)\nu - Q(\nu)b\bigr) = 2(\nu\cdot b)\,I\nu$, so
\[
y^{2} = \psi\nu(\rev{\psi}\psi)\nu\rev{\psi} = \lambda\,I\,y, \qquad \lambda = 2\,\nu\cdot b .
\]
With $y = j + sI$ and $I$ central, $y^{2} = Q(j) + s^{2}I^{2} + 2sIj$ and $\lambda Iy = \lambda Ij + \lambda sI^{2}$. The grade-$4$ parts give $(2s - \lambda)j = 0$ and the scalar parts $Q(j) + s^{2}I^{2} = \lambda sI^{2}$. Where $j \ne 0$ this gives $\lambda = 2s$ and $Q(j) = s^{2}I^{2}$; $j$ is a nonzero polynomial map (check), so $j \ne 0$ on a dense open set of $S$, and the identity $Q(j) = s^{2}I^{2}$ between polynomials in the coordinates of $\psi$ then holds on all of $S$. In both towers $Q(e_{t})I^{2} = 1$ ($-1$ times $-1$ in $\R^{4,1}$, $1$ times $1$ in $\R^{1,4}$), so $m^{2}(j) = Q(e_{t})Q(j) = s^{2}$. The check confirms $y^{2} = 2sIy$, $Q(e_{t})Q(j) = s^{2}$ and $\langle xx^{\dagger}\rangle_{0} = 2s^{2}$ for $x = \rev{\psi}\psi$ as polynomial identities in the eight coordinates of $\psi$, in both towers; the last is the agreement with Theorem \claimref{dark-tower-matter.where-mass-begins}(ii), $m^{2}(j(\psi)) = \tfrac12\langle xx^{\dagger}\rangle_{0}$.

(iv) The same linear system on $S \oplus S$, the bivectors acting on each summand, has a solution space of dimension $3$ in $\R^{2,1}$ and $\R^{1,2}$ and $4$ in $\R^{1,3}$ and $\R^{3,1}$ (exact rank). The maps $j(\psi_{1})$ and $j(\psi_{2})$ are two independent solutions and their sum is the total momentum of the Definition (the check confirms that the sum solves the system); the remaining solutions are cross terms, bilinear in $\psi_{1}$ and $\psi_{2}$. So covariance does not single out the sum.

(v) By (i) a nonzero spin-covariant quadratic vector is $cj$ with $c \ne 0$ real. Replacing $j$ by $cj$ multiplies every total momentum by $c$ and every squared mass $m^{2} = Q(e_{t})Q(P)$ by $c^{2}$; for $c < 0$ it reflects the half-space of momenta. The stored results that assume H1 use the number of constituents at which mass begins and the groups preserving momentum, which do not change under $j \mapsto cj$ (a reading of their statements, not an audit of their proofs). The remaining assumptions listed in the statement are those the argument does not touch: the existence of the spinor part, quadraticity, covariance, and the additivity that (iv) shows is not forced.
